% ============================================================
%  sections/appendix.tex  —  Supplementary Material (Appendix)
% ============================================================
\appendix

\section{Weather-to-Friction Model Details}
\label{app:friction_model}

This appendix provides the full mathematical formulation of the tire--pavement friction model used in SceneFactory (summarized in \S\ref{sec:weather_friction}).
We implement the modified Average Lumped LuGre (ALL) model of Zhao et al.~\cite{zhao_tire-pavement_2024}, which extends the LuGre bristle-deflection framework to account for pavement texture and water film thickness.

\paragraph{Bristle dynamics.}
The model evolves the average bristle deflection $\bar{z}$ via the ODE
\begin{equation}
\label{eq:app_bristle_ode}
\frac{d\bar{z}}{dt} = v_r - \theta \, Y_R \left[\frac{\sigma_0 \lvert v_r \rvert}{g(v_r)}\,\bar{z} + K \lvert \omega_r \rvert\,\bar{z}\right],
\end{equation}
where $v_r$ is the relative (slip) speed between the tire contact patch and the road, $\omega_r$ is the wheel circumferential speed, $\sigma_0$ is the bristle stiffness, $K = 7/(6L)$ is the ALL approximation constant with $L$ the contact patch length, $\theta$ is a texture influence coefficient (per road surface type), and $Y_R \in [0,1]$ is the contact patch length ratio under hydrodynamic lift.

\paragraph{Stribeck function.}
The steady-state function $g(v_r) = \mu_c + (\mu_s - \mu_c)\exp\bigl(-\lvert v_r / v_s \rvert^\alpha\bigr)$ captures the transition from static friction $\mu_s$ to Coulomb friction $\mu_c$.
The Stribeck speed depends on the water film thickness $h_w$ (in meters) via
\begin{equation}
\label{eq:app_stribeck_speed}
v_s(h_w) = b_1 \exp(1000\,b_2\,h_w + b_3) + b_4,
\end{equation}
where $b_1 = 4.8916$, $b_2 = -7.91$, $b_3 = 3.01$, and $b_4 = 3.40$ are calibrated constants from~\cite{zhao_tire-pavement_2024}.
Because $b_2 < 0$, $v_s$ decreases with water film thickness: thicker water films cause the friction drop-off to onset at lower slip speeds.

\paragraph{Effective friction coefficient.}
Given the steady-state solution $\bar{z}^*$, the effective friction coefficient is
\begin{equation}
\label{eq:app_mu_all}
\mu = \max\!\bigl(\theta\,Y_R - Y_F,\; 0\bigr) \cdot \bigl(\theta\,Y_R\,\sigma_0\,\bar{z}^* + \sigma_1\,\dot{\bar{z}} + \sigma_2\,v_r\bigr),
\end{equation}
where $Y_F \in [0,1]$ is the hydrodynamic force ratio and $\sigma_1$, $\sigma_2$ are viscous damping coefficients.
Under the default calibration from~\cite{zhao_tire-pavement_2024}, both $\sigma_1 = 0$ and $\sigma_2 = 0$, so the effective formula reduces to $\mu = \max(\theta\,Y_R - Y_F,\,0)\cdot\theta\,Y_R\,\sigma_0\,\bar{z}^*$.
Both $Y_R$ and $Y_F$ depend on vehicle speed, water film thickness, tire geometry (contact patch dimensions, tread radius, wheel radius), and physical constants (water density and viscosity).
When the contact factor $\theta\,Y_R - Y_F$ drops to zero, the model predicts full hydroplaning.
Our implementation numerically integrates Equation~\ref{eq:app_bristle_ode} to steady state using exponential integration ($\Delta t = 10^{-3}$\,s, convergence threshold $\epsilon = 10^{-7}$), matching the paper's calibrated parameters (Table~4 and Table~5 in~\cite{zhao_tire-pavement_2024}).

\paragraph{Road surface presets.}
The texture influence coefficient $\theta$ and the texture amplitude $T_d$ vary by pavement type.
We provide three presets calibrated to the values in~\cite{zhao_tire-pavement_2024}:
\emph{Asphalt Concrete} (AC, $\theta{=}1.00$, $T_d{=}0.65$\,mm),
\emph{Stone Mastic Asphalt} (SMA, $\theta{=}1.09$, $T_d{=}0.80$\,mm), and
\emph{Open-Graded Friction Course} (OGFC, $\theta{=}1.21$, $T_d{=}1.08$\,mm).
OGFC, designed for drainage, retains higher friction under wet conditions; AC, the most common pavement, serves as the baseline.

\paragraph{PhysX integration pipeline.}
Each world in the training stage is assigned a friction configuration specifying the road surface type and the water film thickness $h_w$ (in mm).
At environment construction time, the pipeline feeds these values into the modified ALL model (Equation~\ref{eq:app_mu_all}) at a reference speed of 13.89\,m/s ($\approx$50\,km/h) and two slip ratios (0.15 for static, 0.80 for dynamic) to produce per-world static and dynamic friction coefficients $\mu_s$ and $\mu_d$.
The effective coefficient $\mu_{\mathrm{eff}}$ (static by default) is used downstream as follows.

The ground-plane material uses the sysid-tuned friction scale for the dry-asphalt surface type (Section~\ref{sec:vehicle_sysid}), with static friction $\mu^{\mathrm{ground}}_s = \min(1.0,\; f_{\mathrm{surface}} \cdot \sqrt{f_{\mathrm{lon}} \cdot f_{\mathrm{lat}}})$ and dynamic friction $\mu^{\mathrm{ground}}_d = 0.95\,\mu^{\mathrm{ground}}_s$.
PhysX resolves contact friction using the minimum of the ground and tire material values (combine mode \texttt{min}).
Because all cloned environments share a single global ground plane, this material is uniform across worlds.
Per-world friction variation is currently communicated to the policy only through the observation (Equation~\ref{eq:weather_obs}).

\chattodo{\textbf{ACTION ITEM --- FIGURE:} Consider a plot of $\mu_{\mathrm{eff}}$ vs.\ $h_w$ (0--1.0\,mm) for the three surface presets (AC, SMA, OGFC) at $v{=}$13.89\,m/s.}


% ============================================================
\section{Observation Space Details}
\label{app:observations}
% ============================================================

Each agent receives a 1{,}929-dimensional observation vector assembled from four groups.
All spatial quantities are expressed in the agent's body frame via a 2D rotation by $-\psi$ (ego heading), so the policy sees an ego-centric view of the world.

\paragraph{Ego state (7-dim).}
The ego state encodes the agent's goal and velocity in body-frame coordinates:
\begin{enumerate}[leftmargin=*, itemsep=1pt]
  \item Goal position $(g_x^b, g_y^b)$, each divided by the scene bounding-box half-size ($B/2 = 100$\,m).
  \item Heading error to the goal, encoded as $(\sin\Delta\psi,\,\cos\Delta\psi)$.
  \item Euclidean goal distance $\|g^b_{xy}\|$, divided by $B/2$.
  \item Body-frame velocity $(v_x^b, v_y^b)$, each divided by 10\,m/s.
\end{enumerate}
When the weather-to-friction module is active, the 4-dim weather context (Equation~\ref{eq:weather_obs}) is appended to the ego group, yielding an 11-dim input to the ego encoder.

\paragraph{Road geometry context ($K_r \times 5 = 1{,}750$-dim).}
Each agent selects $K_r = 350$ road-segment midpoints from the pre-loaded GPU metadata tensors (\S\ref{sec:data_processing}).
The selection operates in two steps: first, all road points within a radius $r = 10$\,m of the ego position are identified via a squared-distance threshold; second, these candidate points are sorted by their original polyline index to preserve road topology, and the first $K_r$ are taken.
If fewer than $K_r$ points fall within the radius, the remaining slots are zero-padded.
Each selected point contributes five features:
\begin{enumerate}[leftmargin=*, itemsep=1pt]
  \item Relative position $(\Delta x^b, \Delta y^b)$ in ego body frame, each divided by $r$.
  \item Polyline type code, divided by a normalizer ($= 20$).
  \item Road direction unit vector $(d_x^b, d_y^b)$ in ego body frame.
\end{enumerate}

\paragraph{Neighbor vehicles ($K_v \times 7 = 168$-dim).}
The $K_v = 24$ nearest alive agents (by Euclidean distance in world XY) are selected per ego agent.
Self-distance and dead-agent distances are set to infinity before the argsort.
Each neighbor contributes seven features:
\begin{enumerate}[leftmargin=*, itemsep=1pt]
  \item Relative position $(\Delta x^b, \Delta y^b)$ in ego body frame, each divided by $B/2$.
  \item Chassis length and width, each divided by $B/2$.
  \item Relative heading $\Delta\psi$, divided by $\pi$.
  \item Speed $\|v_{xy}\|$, divided by 10\,m/s.
  \item Pairwise time-to-collision (TTC), computed via a swept-circle test.
    Each vehicle is approximated by three equal-radius circles placed along the longitudinal axis at offsets $\{-d,\, 0,\, +d\}$ from the chassis centre, where $r = \max(0.45,\; 0.55\, W)$ and $d = \min\!\bigl(\tfrac{\mathit{WB}}{2},\; \max(0,\; \tfrac{L}{2} - 0.8\,r)\bigr)$.
    For the default vehicle ($L{=}4.0$\,m, $W{=}2.0$\,m, $\mathit{WB}{=}2.6$\,m) this gives $r \approx 1.10$\,m and $d \approx 1.12$\,m.
    The test solves a quadratic closest-approach equation for every pair among the $3 \times 3 = 9$ circle combinations and takes the minimum.
    The resulting TTC is clamped to $[0,\, \tau_{\max}]$ and divided by $\tau_{\max} = 10$\,s.
    Neighbors behind the ego or with no collision trajectory receive $\mathrm{TTC}/\tau_{\max} = 1$.
\end{enumerate}
Slots beyond the number of alive neighbors are zero-padded.


% ============================================================
\section{Network Architecture Details}
\label{app:network}
% ============================================================

The observation groups are processed by a late-fusion actor-critic network with per-modality encoders and max-pool set aggregation, following the design of GPUDrive~\cite{kazemkhani_gpudrive_2025}.
We make two modifications: (i)~we use \emph{masked} max-pool (zero-padded slots are filled with $-\infty$ before the $\max$) rather than a plain pool, and (ii)~we replace their Tanh activations with ELU.
The flat observation is split into its three groups and processed by separate encoder MLPs, one set for the actor and one for the critic (no weight sharing).

\begin{itemize}[leftmargin=*, itemsep=1pt]
  \item \textbf{Ego encoder}: $11 \to 64 \to 64$ (two hidden layers, ELU activations).
  \item \textbf{Road encoder}: $5 \to 96 \to 96$, applied independently to each of the $K_r$ points; the $K_r$ output vectors are aggregated via \emph{masked max-pool}, where zero-padded points are filled with $-\infty$ before the $\max$ operation.
  \item \textbf{Vehicle encoder}: $7 \to 96 \to 96$, applied independently to each of the $K_v$ neighbors; aggregated via masked max-pool in the same way.
\end{itemize}

The three embeddings (64 + 96 + 96 = 256) are concatenated and passed through a shared trunk MLP ($256 \to 128 \to 64$, ELU).
The actor head projects to a 3-dim mean vector $\boldsymbol{\mu}$ (one per action), paired with a state-independent learnable log-standard-deviation $\log\boldsymbol{\sigma}$ to form a diagonal Gaussian policy.
The critic head projects to a scalar state-value estimate.
No empirical observation normalization (running mean/std) is used; all features are manually scaled at assembly time as described in Appendix~\ref{app:observations}.

\chattodo{\textbf{ACTION ITEM --- FIGURE:} Create a network architecture diagram showing the three encoder branches, max-pool aggregation, fusion trunk, and actor/critic heads.}


% ============================================================
\section{Action Space Details}
\label{app:actions}
% ============================================================

Each agent outputs a three-dimensional continuous action $\mathbf{a} = (a_{\mathrm{thr}},\, a_{\mathrm{str}},\, a_{\mathrm{brk}})$.
The policy produces values in $[-1,1]^3$; throttle and brake are rectified to $[0,1]$ before decoding, so a zero-mean Gaussian policy naturally idles.
All three channels are decoded into PhysX \emph{effort targets}; the steering PD loop described in Eq.\ref{eq:steer_pd} is computed in software at every physics substep.

\paragraph{Control frequency.}
Physics runs at $1/\Delta t_{\mathrm{phys}} = 120$\,Hz.
Actions are held constant over a decimation window of 4 substeps, yielding a control frequency of 30\,Hz.
Within each substep the steering PD reads the current joint state, so the effective PD loop rate is 120\,Hz.

\paragraph{Brake sign memory.}
Brake torque must oppose wheel rotation, so its sign is $-\mathrm{sign}(\dot{q}_{\mathrm{wheel}})$.
When $|\dot{q}_{\mathrm{wheel}}| < 10^{-4}$\,rad/s the sign becomes unreliable; we therefore latch the last observed sign and continue applying it until rotation resumes.
This prevents torque chatter near zero velocity.

\paragraph{Actuation parameters.}
Table~\ref{tab:app_action_params} lists the sysid-tuned values used in all experiments.
Drive torque is applied to the two front wheels only (FWD); brake torque is applied to all four wheels with separate front/rear magnitudes.
The sysid procedure additionally fits per-surface friction scales (longitudinal and lateral) that modulate the effective contact friction for each road surface type; these are relevant only for multi-surface scenarios (\S\ref{sec:weather_friction}).

\begin{table}[h]
  \centering\small
  \caption{Action-decoding parameters (sysid-tuned values).}
  \label{tab:app_action_params}
  \begin{tabular}{@{}lll@{}}
    \toprule
    Parameter & Symbol & Value \\
    \midrule
    Max steer angle        & $\theta_{\max}$               & $0.450$\,rad ($\approx 25.8^{\circ}$) \\
    Steering PD $K_p$      & $K_p$                          & $1839.5$\,N$\cdot$m/rad \\
    Steering PD $K_d$      & $K_d$                          & $110.5$\,N$\cdot$m$\cdot$s/rad \\
    Steering effort clamp  & $\tau_{\mathrm{steer,max}}$   & $1200$\,N$\cdot$m \\
    Drive torque           & $\tau_{\mathrm{drive,max}}$   & $600.7$\,N$\cdot$m \\
    Brake torque (front)   & $\tau_{\mathrm{brk,F}}$      & $1090.5$\,N$\cdot$m \\
    Brake torque (rear)    & $\tau_{\mathrm{brk,R}}$      & $980.7$\,N$\cdot$m \\
    Wheel mass             & $m_{\mathrm{wheel}}$          & $37.5$\,kg \\
    Wheel inertia scale    & $s_{\mathrm{inertia}}$        & $1.094$ \\
    Suspension stiffness   & $k_{\mathrm{susp}}$           & $1080.8$\,N/m \\
    Suspension damping      & $c_{\mathrm{susp}}$           & $2764.3$\,N$\cdot$s/m \\
    Yaw stability damping  & $\lambda_{\mathrm{yaw}}$      & $10.6$\,N$\cdot$m$\cdot$s/rad \\
    \bottomrule
  \end{tabular}
\end{table}


% ============================================================
\section{Reward Design Details}
\label{app:rewards}
% ============================================================

The per-step reward is the sum of ten terms: four sparse event rewards that also trigger agent termination, and six dense shaping rewards applied at every active step.
Tuning these weights required extensive iteration; the values reported here correspond to the configuration used in all experiments.

\paragraph{Sparse event rewards.}
Four mutually exclusive events terminate an agent and deliver a one-shot reward:
\begin{enumerate}[leftmargin=*, itemsep=1pt]
  \item \textbf{Goal reach} ($+45.0$).
  Awarded when the ego position is within $d_{\mathrm{goal}} = 3.0$\,m of the goal in XY.

  \item \textbf{Collision} ($-6.0$).
  Triggered when the maximum contact-sensor force exceeds 25\,N.
  A warm-up window of 24 steps after spawn suppresses false positives from initial settling.

  \item \textbf{Crash} ($-10.0$).
  Triggered when the vehicle falls below the ground plane ($z_{\mathrm{rel}} < 0$), drifts more than 100\,m from its start position, or tilts beyond a gravity-vector threshold ($g^b_z > -0.15$, indicating a rollover or severe pitch).

  \item \textbf{Lane-forbidden} ($-20.0$).
  Triggered when any of the vehicle's three bounding circles overlaps a road-edge segment (Waymo types 15--16), detected via an oriented-bounding-box overlap test.
\end{enumerate}
After termination, the agent is teleported off-stage and all subsequent dense rewards are masked to zero.

\paragraph{Dense shaping rewards.}
Six terms provide per-step signal for active agents:
\begin{enumerate}[leftmargin=*, itemsep=1pt]
  \item \textbf{Route progress} (weight $w = 2.0$).
  The displacement between consecutive steps is projected onto the tangent of the nearest driving-lane segment (Waymo types 1--2), oriented toward the goal:
  $r_{\mathrm{prog}} = \mathrm{clamp}\bigl((\mathbf{p}_t - \mathbf{p}_{t-1}) \cdot \hat{\mathbf{t}}_{\mathrm{lane}},\; -2,\; +2\bigr) \cdot w$.
  This rewards forward progress along the route and penalizes backwards motion.

  \item \textbf{Geometric lane-keeping} (weight $w = 0.08$).
  A continuous quality score combines lateral offset and heading alignment with the nearest driving lane:
  \[
    q = \exp\!\bigl(-(\ell / \sigma)^2\bigr) \cdot \bigl[(1 - w_h) + w_h \cdot \max(0,\, \cos(\psi - \psi_{\mathrm{lane}}))\bigr],
  \]
  where $\ell$ is the signed lateral offset, $\sigma = 1.75$\,m is the tolerance, and $w_h = 0.8$ is the heading weight.
  The reward is $r_{\mathrm{lane}} = q \cdot w$.

  \item \textbf{Off-road penalty} (weight $w = -0.5$).
  Applied when the vehicle's lateral offset exceeds 3.25\,m or its distance to the nearest lane exceeds 6.0\,m.

  \item \textbf{Idle penalty} (weight $w = -0.05$).
  Applied when the vehicle speed drops below 1.0\,m/s, discouraging the policy from learning to stop and wait.

  \item \textbf{Inter-vehicle TTC penalty.}
  The minimum pairwise TTC across all neighbors is converted to a penalty:
  \[
    p(\tau) = \min\!\bigl(\alpha_v / \max(\tau,\, \tau_{\mathrm{floor}}),\; p_{\mathrm{max},v}\bigr),
  \]
  with $\alpha_v = 0.10$, $p_{\mathrm{max},v} = 0.35$, and $\tau_{\mathrm{floor}} = 0.5$\,s.
  The reward is $r_{\mathrm{ttc}} = -p(\tau)$.

  \item \textbf{Road-edge TTC penalty.}
  The same $\alpha/\tau$ penalty form is applied to the closest road-edge segment within 40\,m ahead:
  $\alpha_e = 0.07$, $p_{\mathrm{max},e} = 0.40$, $\tau_{\mathrm{floor}} = 0.5$\,s.
  This encourages the policy to slow down when approaching road boundaries.
\end{enumerate}

\begin{table}[t]
  \centering\small
  \caption{Reward components used in all experiments.  Sparse rewards are delivered once and trigger agent termination; dense rewards are applied at every active step.}
  \label{tab:app_reward_terms}
  \begin{tabular}{@{}llrll@{}}
    \toprule
    Term & Type & Weight & Sign & Key threshold \\
    \midrule
    Goal reach         & Sparse & $45.0$  & $+$ & dist $\le 3.0$\,m \\
    Collision          & Sparse & $6.0$   & $-$ & force $\ge 25$\,N, 24-step warmup \\
    Crash              & Sparse & $10.0$  & $-$ & fall / drift $> 100$\,m / tilt \\
    Lane-forbidden     & Sparse & $20.0$  & $-$ & road-edge OBB overlap \\
    \midrule
    Route progress     & Dense  & $2.0$   & $\pm$ & clamp $[-2,+2]$\,m/step \\
    Lane-keeping       & Dense  & $0.08$  & $+$ & $\sigma{=}1.75$\,m, $w_h{=}0.8$ \\
    Off-road           & Dense  & $0.5$   & $-$ & lat $> 3.25$\,m or dist $> 6.0$\,m \\
    Idle               & Dense  & $0.05$  & $-$ & speed $< 1.0$\,m/s \\
    Inter-vehicle TTC  & Dense  & $\alpha{=}0.10$ & $-$ & max $= 0.35$, floor $= 0.5$\,s \\
    Road-edge TTC      & Dense  & $\alpha{=}0.07$ & $-$ & max $= 0.40$, floor $= 0.5$\,s \\
    \bottomrule
  \end{tabular}
\end{table}

\paragraph{Termination and time-out.}
An agent terminates upon any sparse event (goal, collision, crash, or lane-forbidden).
If no event occurs, the episode ends at $T_{\mathrm{max}} = 50$\,s (1{,}500 steps at 30\,Hz).
Terminated agents are teleported to an off-stage location ($x = 1{,}000$\,m) and excluded from subsequent reward and observation computation.


% ============================================================
\section{System Identification Details}
\label{app:sysid}
% ============================================================

The vehicle model has 20 tunable dynamics parameters: drive, brake (front/rear), and steering torques; steering PD gains and limits; suspension stiffness and damping; wheel mass and inertia scale; chassis center-of-mass offset; lateral and yaw damping coefficients; and per-surface friction scales for three surface types (dry, wet, gravel).
We fit these parameters using the cross-entropy method (CEM).

\paragraph{Reference data.}
A PhysX ``teacher'' vehicle (the Isaac Sim built-in PhysX Vehicle model, which includes a validated tire model, drivetrain, and suspension) executes 139 scripted maneuvers on a multi-surface patch track.
The maneuvers span five tiers: 17 longitudinal (throttle sweep 10--100\%, brake sweep, linear ramps), 64 lateral (step-steer and constant-radius at four throttle levels $\times$ four--five steer amplitudes $\times$ left/right), 18 combined (trail-brake at three throttle$\times$steer$\times$direction), 39 frequency-response (sinusoidal steering at multiple throttle/amplitude/frequency settings, plus chirp sweeps), and 1 surface transition across dry asphalt ($\mu_s{=}1.0$), wet asphalt ($\mu_s{=}0.75$), and gravel ($\mu_s{=}0.60$).
The comprehensive suite covers the full operating envelope from 10\% to 100\% throttle, ensuring that the fitted parameters generalize across the vehicle's speed and lateral-force range.
Each rollout records per-frame chassis position, yaw, speed, yaw rate, wheel angular velocities, and steer angles at 60\,Hz.

\paragraph{Loss function.}
Let $x^{\mathrm{ref}}_{k,t}$ denote the teacher's value for channel $k$ at frame $t$, and $x^{\mathrm{stu}}_{k,t}$ the corresponding student value under a candidate parameter vector.
Each rollout has $T$ frames recorded at 60\,Hz.
The per-channel mean squared error is $\mathrm{MSE}_k = \tfrac{1}{T}\sum_{t=1}^{T}(x^{\mathrm{stu}}_{k,t} - x^{\mathrm{ref}}_{k,t})^2$.
The full objective adds two terminal-frame penalties that emphasize end-of-maneuver accuracy:
\begin{equation}
\label{eq:app_sysid_loss}
\mathcal{L} = \sum_{k \in \mathcal{K}} w_k \cdot \mathrm{MSE}_k
  + w_{\mathrm{pos,final}} \cdot \|\mathbf{p}^{\mathrm{stu}}_T - \mathbf{p}^{\mathrm{ref}}_T\|^2
  + w_{\mathrm{spd,final}} \cdot (s^{\mathrm{stu}}_T - s^{\mathrm{ref}}_T)^2,
\end{equation}
where $\mathcal{K} = \{\text{position}_{xy}, \text{yaw}, \text{speed}, \text{yaw rate}, \text{wheel speed}, \text{steer angle}\}$ with weights $w_k = (1.0, 0.4, 0.35, 0.25, 0.05, 0.05)$ respectively.
The vectors $\mathbf{p}^{\mathrm{stu}}_T, \mathbf{p}^{\mathrm{ref}}_T \in \mathbb{R}^2$ are the student and teacher XY positions at the final frame, and $s^{\mathrm{stu}}_T, s^{\mathrm{ref}}_T$ are the corresponding speeds.
The terminal weights are $w_{\mathrm{pos,final}} = 1.5$ and $w_{\mathrm{spd,final}} = 0.75$.

\paragraph{Multi-stage CEM search.}
The CEM search is divided into five stages to manage the 20-dimensional space:
\begin{enumerate}[leftmargin=*, itemsep=1pt]
  \item \emph{Longitudinal}: tunes drive/brake torques and wheel parameters on the straight-line rollout.
  \item \emph{Steering}: tunes steering gains, suspension, and damping on the turning rollouts.
  \item \emph{Surface}: tunes per-surface friction scales on the transition rollout.
  \item \emph{Refinement}: tunes all 20 parameters jointly on all rollouts with a narrowed search window (18\% of the full range).
  \item \emph{Brake preservation}: re-tunes longitudinal parameters with a tight window (10\%) to prevent regression.
\end{enumerate}
Each stage inherits the best configuration from the previous stage.
The CEM uses a population of 24 candidates, an elite fraction of 25\%, an initial standard deviation of 25\% of each parameter's range, and a minimum standard deviation of 5\% to prevent premature convergence.
A total of 320 trials are allocated across the five stages with weights (30\%, 20\%, 15\%, 20\%, 15\%).

\paragraph{Vehicle-intrinsic vs.\ environment parameters.}
The 20 sysid parameters are \emph{vehicle-intrinsic}: they characterize how the articulation converts actuator commands into motion (drive torque, steering response, suspension compliance, etc.).
They do not change when the road surface changes.
Environment-level friction variation is handled orthogonally by the PhysX contact solver: each environment's wheel colliders are assigned a friction coefficient $\mu$ via their physics material, and the contact solver computes friction forces at the wheel--ground contact points accordingly (using combine mode \texttt{min}).
This separation is deliberate: sysid produces a realistic vehicle model once, and the physics engine then automatically adapts that vehicle's behavior to any friction environment.

\chattodo{\textbf{ACTION ITEM --- FIGURE:} Consider a trajectory overlay figure showing teacher vs.\ student vehicle on one or two maneuvers.}


% ============================================================
\section{Simulator Limitations and Sharp Edges}
\label{app:sharp_edges}
% ============================================================

Transparency about simulator assumptions is essential for interpreting the experimental results.
Below we catalogue the most consequential simplifications, known failure modes, and practical workarounds in SceneFactory, organized by subsystem.

\paragraph{Z-flattening.}
The Waymo Open Motion Dataset provides full 3D polyline coordinates, including road elevation changes across bridges, ramps, overpasses, and multi-level interchanges.
SceneFactory collapses all road geometry to a constant ground height ($z{=}0$), because vehicles drive on a single flat PhysX ground plane per world.
This design decision has two consequences.
First, scenes with elevated structures (bridges, highway overpasses) produce visually overlapping road segments at ground level, which can confuse the $k$-nearest-road-point observation by presenting the agent with road points from an overhead or underpass road that it cannot physically reach.
Second, ramp grades and banked curves are lost, removing a class of driving challenges (hill starts, crest reactions) that affect real-world vehicle dynamics.
We mitigate the first issue through scene curation and by filtering road points whose original Waymo $z$ coordinate deviates significantly from the ground plane, but the fundamental limitation remains: SceneFactory is a 2.5D simulator operating on 3D data.

\paragraph{Scene curation.}
Not all Waymo scenarios are suitable for flat-ground training.
We developed an interactive \emph{Config Wizard} tool (\texttt{scene\_factory\_config\_wizard.py}) that presents each candidate scene as a bird's-eye plot showing lane-center polylines (Waymo types 1--2, blue), road-edge polylines (types 15--16, red), and valid start--goal pairs connected by origin--destination lines.
A human reviewer accepts or rejects each scene via keyboard input (y/n/b/q), with undo support.
The tool supports both a lightweight Matplotlib mode for rapid triage and an Isaac Sim preview mode that renders the full USD road geometry with 3D start/goal markers.

From a pool of ${\sim}100$ converted Waymo scenarios, we curated 64 training scenes and 64 held-out evaluation scenes, rejecting 23 training candidates and 49 evaluation candidates.
Common rejection reasons include: (i)~multi-level interchanges where z-flattening produces overlapping geometry that would poison the road-point observation, (ii)~scenes where origin--destination pairs straddle a road-edge polyline, making the lane-forbidden penalty unavoidable, and (iii)~scenes with fewer than two valid controllable agents after spawn filtering.
The wizard produces a complete training configuration (scene list, environment count, reward weights, and curriculum stage) from scratch, so a researcher can go from raw Waymo data to a reproducible training run with minimal manual editing.

\paragraph{Waymo road-edge labeling.}
SceneFactory treats Waymo polyline types 15 and 16 as forbidden road edges for both the lane-forbidden termination event and the road-edge TTC penalty.
In practice, these labels are not uniformly reliable across the dataset.
As noted by Kazemkhani et al.~\cite{kazemkhani_gpudrive_2025}, the Waymo dataset labels parking-lot entrances as road edges.
Some tracked vehicles in these scenarios have their origin on the main road and their destination inside the parking area (or vice versa), so the recorded trajectory necessarily crosses a road-edge polyline.
When such a scenario is used for training, the agent faces an impossible conflict: reaching the goal requires crossing a boundary that triggers the lane-forbidden penalty ($-20.0$) and immediate termination.
These cases inject contradictory supervision into the training signal, degrading both goal-reach rate and lane-safety metrics.
Other scenarios lack road-edge annotations entirely where a physical curb or barrier clearly exists, creating blind spots in the lane-safety reward.
Scene curation removes the worst cases, but residual label noise remains and likely contributes to the modest lane-forbidden violation rates observed in evaluation.

\paragraph{Moving-only agent assumption.}
The spawn filter retains only Waymo agents whose start and goal positions are separated by more than 3.0\,m (the goal-reached threshold).
Agents below this threshold are treated as ``already at goal'' and excluded entirely.
Because the Waymo dataset records start and end positions from the logged trajectory, vehicles that remain parked, idle, or nearly stationary throughout the 9.1\,s recording window have near-zero start--goal displacement and are filtered out.
As a result, \emph{parked vehicles do not exist in the simulation}: a scene in which a moving vehicle must navigate past a row of curbside-parked cars will contain only the moving vehicle.
This simplification removes an entire class of realistic interactions (door-zone avoidance, double-parked obstacle negotiation, parallel parking maneuvers) and inflates the effective lane width available to the agent.
A future extension could retain filtered agents as static rigid-body obstacles rather than discarding them.

\paragraph{PointInstancer instability under Fabric.}
Road segments are rendered as USD \texttt{PointInstancer} prims for compactness (one instancer per road type per world, with per-instance position, orientation, and scale arrays).
We discovered that PointInstancer prims interact poorly with Isaac Lab's Fabric state management layer: visibility updates and prototype-table initialization are unreliable, producing intermittent rendering artifacts (invisible or misplaced road segments) that do not affect physics but corrupt visual debugging and video captures.
An explicit-prim fallback mode is provided (one \texttt{UsdGeom.Cube} per road segment) for diagnostic and stage-export workflows.
Production training disables Fabric entirely (\texttt{use\_fabric: false}) to avoid related issues with road metadata access, at a modest throughput cost.

\paragraph{Collision-sensor warmup.}
When vehicles are spawned or teleported to their start poses, PhysX requires several substeps to settle the suspension and establish stable wheel--ground contact.
During this transient, contact sensors report large spurious forces as the chassis drops onto the ground plane.
Without mitigation, these forces trigger false collision terminations at the start of every episode.
We suppress collision detection for the first 24 environment steps (${\approx}0.8$\,s) after each spawn, which is sufficient for the suspension to settle.
This warmup window introduces a brief period at episode start during which actual inter-vehicle collisions are unpenalized, though in practice vehicles are spawned with sufficient separation that early collisions are rare.

\paragraph{Teleport-only reset.}
Isaac Lab's default environment reset tears down and reconstructs PhysX actors, which is prohibitively expensive for multi-agent scenes with per-environment road geometry.
SceneFactory uses a \emph{teleport-only} reset strategy: after a single full initialization reset, subsequent episode resets reposition vehicles by writing directly to root-state tensors (position, orientation, linear and angular velocity) without rebuilding the physics scene.
This is ${\sim}10{\times}$ faster but has a subtle limitation: accumulated PhysX internal state (joint impulse caches, contact-pair broadphase data) is not fully cleared.
In rare cases this causes a vehicle to retain a residual velocity or contact state from the previous episode for one or two substeps after teleport.
The collision warmup window (above) masks most such artifacts.

\paragraph{Uniform friction surface.}
Per-world friction variation is implemented by writing to each vehicle's wheel collider physics material at environment construction time.
Because the global ground plane has friction $\mu{=}1.0$ and PhysX uses combine mode \texttt{min}, the effective contact friction is $\min(\mu_{\mathrm{wheel}}, 1.0) = \mu_{\mathrm{wheel}}$, giving per-environment control.
However, friction remains spatially uniform within each world: every wheel contact with the ground plane uses the same coefficient.
Real roads exhibit local friction variation (wet patches, oil spills, transitions between asphalt and concrete), and scenarios such as a slippery lane adjacent to a dry lane cannot be represented.
Modeling local friction would require either a segmented ground mesh with per-face materials or runtime material-property writes keyed to vehicle position, neither of which is straightforward under PhysX~5's current contact pipeline.

\paragraph{PhysX substep granularity.}
Physics runs at 120\,Hz with a decimation factor of 4, yielding 30\,Hz control.
The 120\,Hz rate is a compromise between simulation fidelity and throughput: higher rates (240\,Hz or 480\,Hz) would improve suspension stability and contact resolution but proportionally reduce training speed.
Our sysid calibration (\S\ref{sec:vehicle_sysid}) was performed at this rate, so the fitted parameters absorb any 120\,Hz-specific artifacts; re-calibration would be needed if the substep rate were changed.

\paragraph{Clone-in-Fabric disabled.}
Isaac Lab's \texttt{clone\_in\_fabric} optimization, which clones environment state entirely in the Fabric layer for maximum GPU throughput, is disabled in SceneFactory because it is incompatible with the post-clone road injection workflow described in \S\ref{sec:data_processing}.
Fabric assumes that all environments are structurally identical after cloning; SceneFactory's per-environment road prims violate this assumption, causing metadata reads to return stale or incorrect values.
Disabling Fabric increases CPU--GPU synchronization overhead but is necessary for correctness.

\paragraph{Constant-velocity TTC.}
Both the inter-vehicle and road-edge TTC computations assume constant velocity over the prediction horizon.
This systematically overestimates collision risk during deceleration (the vehicle is predicted to continue at its current speed toward the obstacle) and underestimates risk during acceleration.
A quadratic (constant-acceleration) TTC formulation would be more accurate, particularly in approach zones where vehicles are actively braking, but would increase the per-step computational cost of the $K_v{\times}9$ pairwise circle tests.


% ============================================================
\section{Throughput Analysis Details}
\label{app:throughput}
% ============================================================

\paragraph{Timing breakdown.}
Table~\ref{tab:app_timing_breakdown} reports the per-step timing breakdown for both the Baseline and SceneFactory pipelines at 64 worlds (1{,}024 agent slots, all active), from the Experiment~A scaling benchmark on scene \texttt{000077}.

On the Baseline side, the PhysX rigid-body solver accounts for roughly half of wall-clock time (2{,}934\,ms, 51.6\%); this cost is irreducible and shared by both pipelines.
The observation builder consumes another 1{,}456\,ms (25.6\%), because it iterates over each of the 1{,}024 active agents in a Python loop, issuing per-prim USD state queries and computing ego-frame transforms with scalar arithmetic.
Geometric lane features, TTC computation, and action application add a further 22\% through similar per-agent loops.
The total per-step cost is \textbf{5{,}683\,ms}.

On the SceneFactory side, the same 64-world configuration completes an environment step in \textbf{135\,ms}---a \textbf{42$\times$} reduction.
Physics remains the largest single cost (53.8\,ms), but the GPU-batched solver is 55$\times$ faster than the Baseline's single-scene PhysX dispatch.
Observations, rewards, dones, and actions collectively take only 78\,ms because they execute as batched \texttt{torch} tensor operations with no per-agent Python loop.

\begin{table}[t]
\centering
\caption{Per-step timing breakdown at 64 worlds (1{,}024 agent slots, all active) from the Experiment~A benchmark on scene \texttt{000077}.
Baseline sub-phases are grouped into the matching SceneFactory category.
SceneFactory sub-phases may overlap on the GPU, so they sum to more than the measured wall-clock total; the \textbf{Total} row reports the true measured time.}
\label{tab:app_timing_breakdown}
\small
\begin{tabular}{lrrr}
\toprule
Phase & Baseline (ms) & SceneFactory (ms) & Speedup \\
\midrule
Physics solver         & 2{,}934 &  53.8 & 55$\times$ \\
Observation builder    & 1{,}772 &  23.3 & 76$\times$ \\
Reward \& termination  &     629 &  29.9 & 21$\times$ \\
Action application     &     299 &  24.9 & 12$\times$ \\
Reset / overhead       &      49 &  32.4 &  1.5$\times$ \\
\midrule
\textbf{Total (wall-clock)} & \textbf{5{,}683} & \textbf{135} & \textbf{42$\times$} \\
\bottomrule
\end{tabular}
\\[4pt]
{\footnotesize
\emph{Baseline grouping:}
Observation builder = obs.\ construction (1{,}456\,ms) + geometric lane features (316\,ms);
Reward \& termination = road-edge TTC (236) + vehicle TTC (191) + TTC state queries (136) + contact scan (66).
\emph{SceneFactory grouping:}
Reward \& termination = reward computation (18.4) + done/termination (11.5);
Physics solver = write (33.0) + simulate (15.9) + read-back (4.9).}
\end{table}

\paragraph{Why the Baseline is slow.}
SceneFactory eliminates Baseline bottlenecks through three architectural changes:
\begin{enumerate}[leftmargin=*, itemsep=2pt]
  \item \textbf{Batched joint-state tensors.}
        The custom URDF articulated vehicle (\S\ref{sec:vehicle_sysid}) exposes positions, velocities, and orientations as GPU-resident tensors of shape $(\texttt{num\_envs}, \texttt{num\_agents}, d)$ via Isaac Lab's \texttt{Articulation} API.
        Observation construction reduces to a handful of batched \texttt{torch} operations with no Python per-agent loop.
  \item \textbf{Cloned PhysX scenes.}
        Isaac Lab's environment cloning creates $N$ independent PhysX solver islands, enabling the GPU broad-phase and narrow-phase to dispatch them in parallel.
        The Baseline places all worlds in a single PhysX scene, limiting solver parallelism.
  \item \textbf{GPU-resident training loop.}
        RSL-RL reads observations, rewards, and dones as CUDA tensors directly from the environment; no CPU round-trip occurs during training.
        The Baseline converts observations to NumPy for Stable-Baselines3's rollout buffer, incurring $O(\text{agents} \times \text{obs\_dim})$ bytes of device-to-host transfer every step.
\end{enumerate}


% ============================================================
\section{PPO Hyperparameters}
\label{app:ppo_hparams}
% ============================================================

Table~\ref{tab:app_ppo_hparams} lists the full PPO configuration used in all experiments.
The learning rate follows an adaptive schedule that halves or doubles the rate to maintain the per-update KL divergence near a target of 0.01.

\begin{table}[h]
\centering
\caption{PPO hyperparameters.}
\label{tab:app_ppo_hparams}
\small
\begin{tabular}{ll}
\toprule
Parameter & Value \\
\midrule
Rollout horizon           & 128 steps \\
Mini-batches              & 16 \\
Learning epochs           & 5 \\
Learning rate (initial)   & $2.09 \times 10^{-4}$ \\
LR schedule               & adaptive (KL target 0.01) \\
Discount $\gamma$         & 0.99 \\
GAE $\lambda$             & 0.98 \\
Clip $\epsilon$           & 0.15 \\
Entropy coefficient       & $3 \times 10^{-5}$ \\
Value loss coefficient    & 1.0 \\
Max gradient norm         & 1.0 \\
\bottomrule
\end{tabular}
\end{table}
